\documentclass[12pt]{article}
\usepackage[utf8]{inputenc}
\usepackage[russian]{babel}
\usepackage{amsmath}
\usepackage{geometry}
\geometry{a4paper, margin=2.5cm}

\title{Реферат: три задачи интеллектуальной обработки данных}
\author{}
\date{}

\begin{document}
\maketitle

\section{Описание системы}

Рассматриваются три независимые задачи обработки данных, объединённые общим
пайплайном <<признаки $\to$ модель $\to$ оценка качества>>:

\begin{enumerate}
  \item \textbf{Детекция и подсчёт объектов на изображении} — система получает
        на вход растровое изображение со множеством объектов (монеты),
        выполняет бинаризацию и сегментацию и возвращает их число и положение.
  \item \textbf{Классификация тональности отзывов} — система получает текст
        отзыва покупателя о магазине/ресторане и определяет его тональность
        (положительный/отрицательный).
  \item \textbf{Диагностика неисправностей по сигналам датчиков} — система
        получает временные ряды вибрации и тока электродвигателя и определяет
        техническое состояние: норма или один из видов неисправности
        (дисбаланс, дефект подшипника, несоосность).
\end{enumerate}

Во всех трёх случаях задача сводится к отображению входных данных $x$ в метку
класса $y$ из конечного множества состояний.

\section{Математические формулировки задач}

\subsection{Детекция и подсчёт объектов}

Изображение $I(u,v) \in [0,1]$ бинаризуется порогом Отсу $t^*$, минимизирующим
внутриклассовую дисперсию:
\begin{equation}
  t^* = \arg\min_t \; \omega_0(t)\sigma_0^2(t) + \omega_1(t)\sigma_1^2(t),
\end{equation}
где $\omega_i, \sigma_i^2$ — вес и дисперсия класса фон/объект. Бинарная маска
$B(u,v) = \mathbb{1}[I(u,v) < t^*]$ далее разбивается на связные компоненты
watershed-преобразованием по карте расстояний $D(u,v)$, а число объектов
$N = |\{R_1,\dots,R_N\}|$ равно числу найденных областей $R_i$ с площадью
$|R_i| \ge S_{\min}$.

\subsection{Классификация тональности отзывов}

Текст отзыва $d$ представляется TF-IDF вектором $x \in \mathbb{R}^p$:
\begin{equation}
  x_j = \mathrm{tf}_j(d) \cdot \log\frac{N_{doc}}{1+\mathrm{df}_j},
\end{equation}
где $\mathrm{tf}_j$ — частота $j$-й n-граммы в $d$, $\mathrm{df}_j$ — число
документов, содержащих эту n-грамму. Классификатор оценивает
\begin{equation}
  P(y=1\mid x) = \sigma(w^\top x + b), \qquad \sigma(z) = \frac{1}{1+e^{-z}},
\end{equation}
а параметры $w, b$ находятся минимизацией логистической функции потерь
$\mathcal{L} = -\sum_i \big[y_i\log \hat y_i + (1-y_i)\log(1-\hat y_i)\big]$.

\subsection{Диагностика неисправностей}

Сигнал $s(t)$ (вибрация/ток) описывается вектором признаков
$x = (x_{\mathrm{rms}}, x_{\mathrm{crest}}, x_{f_0}, x_{\mathrm{hf}}, \dots)$, где
\begin{equation}
  x_{\mathrm{rms}} = \sqrt{\frac{1}{T}\int_0^T s^2(t)\,dt}, \qquad
  x_{\mathrm{hf}} = \frac{\sum_{f>f_c} |S(f)|^2}{\sum_{f} |S(f)|^2},
\end{equation}
$S(f)$ — амплитудный спектр Фурье сигнала. Задача — многоклассовая
классификация $\hat y = \arg\max_{k} P(y=k\mid x)$, $k\in\{$норма, дисбаланс,
подшипник, несоосность$\}$.

\section{Нейросетевые решения}

Описанные выше модели (логистическая регрессия, случайный лес) являются
базовыми. Все три задачи естественно обобщаются на нейросетевые архитектуры:

\begin{itemize}
  \item \textbf{Детекция объектов:} свёрточная сеть (CNN) типа U-Net для
        семантической сегментации масок объектов или детектор YOLO/Faster R-CNN
        для прямого предсказания ограничивающих рамок и их числа $N$.
  \item \textbf{Тональность отзывов:} вместо TF-IDF используется
        полносвязный/рекуррентный слой над эмбеддингами слов, либо
        предобученный трансформер (BERT), дающий $P(y\mid x) = \mathrm{softmax}(f_\theta(x))$,
        где $f_\theta$ — нейросеть с обучаемыми весами $\theta$.
  \item \textbf{Диагностика неисправностей:} одномерная свёрточная сеть
        (1D-CNN) или LSTM, принимающая сырой временной ряд $s(t)$ без
        ручного извлечения признаков и обучаемая end-to-end на
        многоклассовую кросс-энтропию
        $\mathcal{L} = -\sum_k y_k \log \hat y_k$.
\end{itemize}

Во всех случаях переход к нейросетевому решению заменяет этап ручного
конструирования признаков обучаемым представлением, что потенциально
повышает точность при достаточном объёме размеченных данных.

\end{document}
